\documentclass[11pt,a4paper]{article}

\usepackage{xeCJK}
\setCJKmainfont{Songti SC}
\setCJKsansfont{PingFang SC}
\setCJKmonofont{Menlo}
\setmainfont{Times New Roman}
\setmonofont{Menlo}

\usepackage[margin=2.2cm]{geometry}
\usepackage{amsmath, amssymb, amsthm, mathtools}
\usepackage{bm}
\usepackage{booktabs, tabularx, array, multirow}
\usepackage{graphicx}
\usepackage{xcolor}
\usepackage{hyperref}
\usepackage{enumitem}
\usepackage{caption}
\usepackage{fancyhdr}
\usepackage{tcolorbox}

\hypersetup{
  colorlinks=true,
  linkcolor=blue!60!black,
  urlcolor=blue!60!black,
  citecolor=blue!60!black,
  pdftitle={硬判决 RS+BCH 级联课题——需求解读与技术追问},
  pdfauthor={陈诗阳}
}

\pagestyle{fancy}
\fancyhf{}
\rhead{\footnotesize 硬判决 RS+BCH 级联}
\lhead{\footnotesize 需求解读与技术追问}
\cfoot{\thepage}

\renewcommand{\baselinestretch}{1.2}

\definecolor{accent}{RGB}{30,80,160}
\definecolor{accent2}{RGB}{160,40,40}
\definecolor{accent3}{RGB}{40,120,60}

\newtcolorbox{keybox}[1][]{colback=blue!5, colframe=accent, boxrule=0.4mm,
  arc=1mm, left=2mm, right=2mm, top=1mm, bottom=1mm, #1}
\newtcolorbox{riskbox}[1][]{colback=orange!5, colframe=orange!70!black, boxrule=0.4mm,
  arc=1mm, left=2mm, right=2mm, top=1mm, bottom=1mm, #1}
\newtcolorbox{actionbox}[1][]{colback=green!5, colframe=accent3, boxrule=0.4mm,
  arc=1mm, left=2mm, right=2mm, top=1mm, bottom=1mm, #1}
\newtcolorbox{quotebox}[1][]{colback=gray!8, colframe=gray!50, boxrule=0.3mm,
  arc=1mm, left=2mm, right=2mm, top=1mm, bottom=1mm, #1}

\begin{document}

% ============ Cover ============
\begin{center}
{\Large\bfseries 硬判决 RS+BCH 级联码课题}\\[4pt]
{\large 新需求解读与需要对齐的技术问题}\\[16pt]
{\normalsize 陈诗阳\\
基于 Lagendijk et al.\ 2026 硬件实现论文\\
\vspace{2pt}
\texttt{https://github.com/a-yeyang/efficient-osd-bch-no-ge}}\\[10pt]
\rule{0.9\textwidth}{0.4pt}\\[6pt]
{\footnotesize 2026-07-23}\\
\rule{0.9\textwidth}{0.4pt}
\end{center}

\vspace{6pt}
\begin{keybox}[title=\textbf{文档目的}]
本文档针对老师最新的\textbf{硬判决级联码}任务需求，整理出：
\begin{enumerate}[leftmargin=1.5em, itemsep=1pt, topsep=2pt]
\item Lagendijk 2026 论文的核心技术要点（\S1）
\item 老师新需求的逐条解读（\S2）
\item 相对第一阶段课题的根本变化（\S3）
\item 8 个需要对齐的技术问题（\S4）
\item 建议的执行路线与保守推进策略（\S5）
\end{enumerate}
\end{keybox}

% ============ Section 1: Paper ============
\section{Lagendijk 2026 论文技术要点}

\subsection{论文定位}

\textbf{论文}：Lagendijk et al., ``High-throughput Low-latency Hardware Implementation of BCH Decoders,''
arxiv 2606.17837v1 (EPFL + Eindhoven Univ.\ of Technology, 2026 年 6 月)。

这是一篇\textbf{硬件实现论文}，为 BCH 译码器设计 FPGA/ASIC 架构。核心贡献是对比两种译码路径：

\begin{table}[!ht]
\centering
\begin{tabularx}{\textwidth}{lXXX}
\toprule
\textbf{路径} & \textbf{算法} & \textbf{支持} $t$ & \textbf{特点} \\
\midrule
Conventional & BM (Berlekamp--Massey) + Chien search & 任意 $t$ & 迭代求 ELP + 穷举根 \\
Direct       & 直接求 ELP 的根 (闭式代数变换 + LUT) & $t \le 4$ & 一次性给出根 \\
\bottomrule
\end{tabularx}
\end{table}

\subsection{Direct 译码的关键技术}

对于 $t = 2, 3, 4$，错误定位多项式 (ELP) 的\textbf{根}可以用\textbf{闭式代数变换}直接求出：

\begin{itemize}[leftmargin=1.5em, itemsep=1pt]
\item $t = 2$: ELP 是 2 次多项式 $\Lambda(x) = x^2 + S_1 x + (S_1^3 + S_3)/S_1$，
  变换成 $A(x) = x^2 + x + k$ 标准形式，用\textbf{一个大小 $2^m - 1$ 的 LUT} 直接查根
\item $t = 3$: 变换成 $B(x) = x^3 + x + k$ 或 $C(x) = x^3 + k$，用另一个 LUT
\item $t = 4$: 4 次多项式，分解为两个 2 次多项式的复合，再查表
\end{itemize}

\textbf{硬件收益}：Direct 译码把整个 BCH 译码压缩成\textbf{3--8 个时钟周期}，
而 Conventional BM + Chien 需要 $2t + 2$ 时钟周期 + 全 $n$ 位 Chien 搜索。

\subsection{论文关键结果 (n=256, Table VI)}

\begin{table}[!ht]
\centering
\begin{tabularx}{\textwidth}{llrrrl}
\toprule
Decoder & $t$ & Latency (cyc) & Latency (ns) & FPGA & 面积改善 \\
\midrule
Conv     & 2   & 8             & 16           & XCZU48DR & baseline \\
\textbf{Direct}   & \textbf{2}   & \textbf{3}    & \textbf{6}   & XCZU48DR & \textbf{-62\% 时延, -65\% LUT, +35\% clock speed} \\
Conv     & 3   & 8             & 28           & XCZU48DR & baseline \\
\textbf{Direct}   & \textbf{3}   & \textbf{4}    & \textbf{8}   & XCZU48DR & \textbf{-71\% 时延, -75\% LUT} \\
\bottomrule
\end{tabularx}
\caption{Conventional vs Direct BCH 译码（Lagendijk 2026 Table VI）}
\end{table}

\begin{keybox}[title=\textbf{核心观察}]
Direct $t=2$ 译码相对 Conventional 是\textbf{3/8 = 37.5\% 时延}
（相当于 \textbf{2--3$\times$ 的硬件层面时延压缩}）。这是硬件综合数据，不是软件仿真估算。
\end{keybox}

% ============ Section 2: Advisor's new task ============
\section{老师新需求逐条解读}

\subsection{原文}

\begin{quotebox}
用类似共享拉格朗日插值理念以及这篇文章 (Lagendijk 2026) 的理念，构建级联码：
\textbf{硬判决 RS + 硬判决 BCH（码长 255/128，$t=2$）}，比较级联码与硬判 RS 的
BER 性能与时延参量，KPI：\textbf{级联码时延不高于硬判 RS 的 10\%}。
\end{quotebox}

\subsection{拆解}

\begin{table}[!ht]
\centering
\begin{tabularx}{\textwidth}{lX}
\toprule
\textbf{原文} & \textbf{精确技术含义} \\
\midrule
共享拉格朗日插值理念 &
沿用上个课题的思路：内外码共享 $\alpha$ 幂表、denominator products、syndrome 单元 \\
Lagendijk 论文的理念 &
内码 BCH 采用 \textbf{Direct root finding + LUT 查表}，而非 BM 迭代 \\
硬判决 RS + 硬判决 BCH &
外码 RS 用 BM + Chien + Forney；内码 BCH 用 Direct root finding。\textbf{全硬判}，无 LLR \\
码长 255/128, $t=2$ &
两组 config，BCH 都是 $t=2$。$n=128$ 严格来说要用 extended BCH \\
KPI 级联时延 $\leq$ 硬判 RS 时延 + 10\% &
在时钟周期或 $\mathbb{F}_{2^m}$ ops 口径下，比较级联链路 vs 纯 RS-BM \\
\bottomrule
\end{tabularx}
\end{table}

% ============ Section 3: Differences from previous phase ============
\section{相对第一阶段课题的根本变化}

上一阶段（软判决级联）与本次任务（硬判决级联）的对比：

\begin{table}[!ht]
\centering
\begin{tabularx}{\textwidth}{lXX}
\toprule
\textbf{维度} & \textbf{第一阶段（软判决）} & \textbf{本次（硬判决）} \\
\midrule
内码 BCH   & LLOSD (软判决，Lagrange + TEP)  & Direct root finding (硬判决，LUT) \\
外码 RS    & BM 硬判 / LCC-BR 软判            & \textbf{只有 BM 硬判} \\
目标基线   & 软判基线 (LCC-BR)               & \textbf{硬判基线 (RS-BM)} \\
KPI        & vs LCC-BR：$\le 10\%$ (已满足)  & vs RS-BM：$\le 10\%$ (需极致优化) \\
编码增益   & $\ge 1$ dB                      & 预期 $0.5$--$1$ dB (硬判天生偏弱) \\
主要优化   & Lagrange 共享 + TEP 早期终止    & \textbf{BCH $t=2$ Direct 译码} \\
时延口径   & 微秒级 ($\mathbb{F}_{2^m}$ ops) & \textbf{纳秒级 (硬件周期)} \\
可复用比例 &  ---                             & \textbf{$\sim$70\% 基础设施} \\
\bottomrule
\end{tabularx}
\end{table}

\subsection{可复用的第一阶段成果}

\begin{itemize}[leftmargin=1.5em, itemsep=1pt]
\item PAM4 modem + AWGN 信道 + bit-LLR（虽然本次只用硬判，但 modem 保持不变）
\item RS 编码器 + BM 硬判解码器
\item Lagrange 缓存层（$\alpha$ 幂表 + denominator products）
\item 级联链路框架 + Monte Carlo 仿真驱动
\item 结果分析和报告脚手架
\end{itemize}

\subsection{本次需要新增的模块}

\begin{itemize}[leftmargin=1.5em, itemsep=1pt]
\item \textbf{BCH $t=2$ Direct root finding} (核心新工作)
\item BCH Conventional BM + Chien 作为对比 baseline
\item 硬件时钟周期时延模型（不再只用软件 ops 计数）
\item 若需要：$n=128$ 的 extended BCH 处理
\end{itemize}

% ============ Section 4: Alignment questions ============
\section{需要对齐的 8 个技术问题}

按重要性从高到低排列。前 3 个问题会\textbf{根本决定后续实现路线}，
建议动手写代码前明确回答。

\subsection{Q1. 内码 BCH 的具体译码器}

\begin{actionbox}[title=\textbf{Q1. Direct 还是 Conventional？还是两者对比？}]
老师说``硬判决 BCH''，Lagendijk 论文里给出\textbf{两种}硬判决方案。请问：

\begin{itemize}[leftmargin=1em, itemsep=1pt]
\item [(a)] 只做 Direct root finding ($t=2$，闭式 + LUT)
\item [(b)] 只做 Conventional (BM + Chien)
\item [(c)] 两个都做，用作论文级别的时延消融对比 \textbf{(推荐)}
\end{itemize}

我建议 (c)：因为论文的核心 claim 就是 Direct vs Conventional 的对比，
且我们的 Lagrange 共享理念与 Direct 的 LUT 化最契合（预计算 $\to$ 查表）。
\end{actionbox}

\subsection{Q2. 时延测量口径}

\begin{actionbox}[title=\textbf{Q2. Clock cycles 还是 F$_{2^m}$ ops？}]
第一阶段用 $\mathbb{F}_{2^m}$ ops 计数，级联相对 RS-BM 慢 12$\times$，Case (a) 不通过。
本次因为 BCH 内码只用 Direct $t=2$ (3 cycles)，主要时延优势体现在\textbf{硬件周期}上。请问：

\begin{itemize}[leftmargin=1em, itemsep=1pt]
\item [(a)] 只报 clock cycles，用 Lagendijk 论文里的具体数字（Direct t=2 = 3 cyc, RS-BM = ？ cyc）
\item [(b)] 同时报 $\mathbb{F}_{2^m}$ ops 和 clock cycles \textbf{(推荐)}
\item [(c)] 我们自建 hardware timing model（$P$ 路并行度参数化）
\end{itemize}

推荐 (b)，能覆盖两种视角，也让老师能横向对比第一阶段结果。
\end{actionbox}

\subsection{Q3. 参数配置的具体码率}

因为 BCH $t=2$ 已固定，码率约束需要重新计算：

\begin{table}[!ht]
\centering
\begin{tabularx}{\textwidth}{XXXX}
\toprule
$n$ & BCH & 最匹配的 RS & 总码率 \\
\midrule
$n=255$ & BCH(255, 239) $t=2$ & RS(255, 239) $t=8$ & $0.937 \times 0.937 = 0.879$ \\
$n=127$ & BCH(127, 113) $t=2$ & RS(127, 125) $t=1$ (太弱) & $0.984 \times 0.890 = 0.876$ \\
$n=127$ (备选) & BCH(127, 113) $t=2$ & RS(127, 113) $t=7$ & $0.890 \times 0.890 = 0.792$ \\
\bottomrule
\end{tabularx}
\caption{$t=2$ BCH 下的 RS 参数选择困境}
\end{table}

\begin{actionbox}[title=\textbf{Q3. $n=127$ 的 RS 参数如何选？}]
\begin{itemize}[leftmargin=1em, itemsep=1pt]
\item [(a)] 严格保持码率 0.88 $\to$ RS $t=1$ (纠错能力太弱, 级联几乎无增益)
\item [(b)] 允许码率 $\sim$0.79 $\to$ RS $t=7$ (强纠错, 但违反 12\% 开销约束) \textbf{(推荐)}
\item [(c)] 用 shortened RS 保持精确码率 0.88 (需要额外处理)
\end{itemize}
\end{actionbox}

\subsection{Q4. $n=128$ 是否放宽为 $n=127$}

\begin{actionbox}[title=\textbf{Q4. $n=128$ 用 extended BCH 还是直接 $n=127$？}]
Lagendijk 论文的 Direct root finding 只对 \textbf{primitive BCH} ($n = 2^m - 1$) 定义。
$n=128$ 需要额外推广，且 Table VI 里只报告了 $n=256$ 的具体数字。请问：

\begin{itemize}[leftmargin=1em, itemsep=1pt]
\item [(a)] 直接改成 $n=127$ (更符合论文范式) \textbf{(推荐)}
\item [(b)] 用 extended BCH ($n=128$)，接受 Direct 译码需要小改
\item [(c)] Shortened BCH (从 $n=255$ shortened)
\end{itemize}
\end{actionbox}

\subsection{Q5. Lagrange 共享的具体位置}

\begin{riskbox}[title=\textbf{Q5. Direct + BM 共享什么？}]
Direct BCH 译码用 \textbf{LUT 查表}（不做 Lagrange 插值），RS-BM 用\textbf{多项式 GF 运算}
（也不直接用 Lagrange 插值）。第一阶段的 Lagrange 共享（$\alpha$ 幂表 + denom products）
在本课题的效用会大幅降低。

老师说的``共享拉格朗日插值理念''，本次可能落在：

\begin{itemize}[leftmargin=1em, itemsep=1pt]
\item \textbf{$\alpha$ 幂表}：BCH syndrome ($S_i = r(\alpha^i)$) 与 RS syndrome 使用相同的 $\alpha$ 幂序列 $\to$ 硬件可复用
\item \textbf{Syndrome 求值单元}：两级共用一个 polynomial evaluator，节省硬件 area
\item \textbf{LUT 表结构}：BCH Direct 的 $\{\}_A$ 表可在编码时预填 $\alpha$ 域信息 (统一 memory 组织)
\end{itemize}

需请老师确认，本次的``Lagrange 共享''是否指这类\textbf{硬件层面的 module 复用}，
而不是软件的 lookup cache（因为软件层面上 Direct 译码基本没有 Lagrange 计算可共享）。
\end{riskbox}

\subsection{Q6. 编码增益预期}

\begin{riskbox}[title=\textbf{Q6. 硬判决级联的编码增益是否达标？}]
硬判 RS + 硬判 BCH 级联，预期编码增益\textbf{比软判弱}。粗略估计：

\begin{itemize}[leftmargin=1em, itemsep=1pt]
\item 纯 RS-BM $n=255, t=8$: FER=$10^{-5}$ 需 $\sim$10.5 dB
\item 级联 (RS-BM $t=8$ + BCH $t=2$): 预期 0.5--1 dB 编码增益 $\to$ FER=$10^{-5}$ 需 $\sim$9.5--10 dB
\end{itemize}

请老师明确：
\begin{itemize}[leftmargin=1em, itemsep=1pt]
\item [(a)] 只要编码增益 $\ge 0.5$ dB 就 OK
\item [(b)] 需要 $\ge 1$ dB, 否则考虑增大 BCH $t$
\item [(c)] 编码增益不是硬 KPI, 重点看时延对比
\end{itemize}
\end{riskbox}

\subsection{Q7. 时延基线的``硬判 RS''具体指哪种}

\begin{actionbox}[title=\textbf{Q7. 硬判 RS 是否使用 Lagendijk 类似的 direct root finding？}]
Lagendijk 论文只给了 BCH 的 direct 译码。但 RS 也可以用类似的 direct 思路，只是复杂度更高。
时延基线（硬判 RS）\textbf{也用 direct root finding} 还是 \textbf{仍用传统 BM + Chien}？

\begin{itemize}[leftmargin=1em, itemsep=1pt]
\item [(a)] 只用 conventional BM + Chien 作基线 \textbf{(推荐, 与论文一致)}
\item [(b)] 也做 direct RS 作对比 (工作量增大)
\end{itemize}
\end{actionbox}

\subsection{Q8. 交付形式}

\begin{actionbox}[title=\textbf{Q8. Python vs Matlab，是否需要 FPGA 综合？}]
Lagendijk 论文最强的部分是 \textbf{FPGA / ASIC 综合报告}（Table V/VI）。请问本次交付：

\begin{itemize}[leftmargin=1em, itemsep=1pt]
\item [(a)] 只做 Python 仿真 + 抽象时延模型 \textbf{(基础交付)}
\item [(b)] 加上 Matlab port
\item [(c)] 加上 FPGA 综合（Verilog + Vivado），产出真实硬件时延数字 \textbf{(工作量翻倍但最有说服力)}
\end{itemize}
\end{actionbox}

% ============ Section 5: Execution plan ============
\section{建议的执行路线}

\subsection{保守推进策略（不等回答的情况下）}

如果需要立即开始，我建议按\textbf{最保守假设}推进，覆盖 Q1--Q8 所有可能路径：

\begin{table}[!ht]
\centering
\begin{tabularx}{\textwidth}{lXl}
\toprule
\textbf{阶段} & \textbf{任务} & \textbf{时长} \\
\midrule
Phase 1 & 实现 BCH $t=2$ Conventional (BM + Chien) & 半天 \\
Phase 2 & 实现 BCH $t=2$ Direct root finding (LUT) & 1 天 \\
Phase 3 & 复用现有 RS-BM 与级联链路 & 半天 \\
Phase 4 & 硬件时延模型 (clock cycles + $\mathbb{F}_{2^m}$ ops 两种口径) & 半天 \\
Phase 5 & 跑 $n=255, t=2$ 完整实验 (Direct + Conv + baseline) & 1 天 \\
Phase 6 & 跑 $n=127$ (若确认可放宽) & 半天 \\
Phase 7 & LaTeX 报告 + 图表 + 时延分析 & 1 天 \\
\bottomrule
\end{tabularx}
\caption{保守推进路线：预计 4.5 天完成到 Python 完整仿真}
\end{table}

\subsection{最小闭环建议}

\begin{keybox}[title=\textbf{推荐先做的最小闭环}]
先做\textbf{$n=255$ + Direct + Conv + RS-BM baseline}，跑通并请老师 review 中间结果，
确认路线正确后再横向扩展到 $n=127$ 和其他配置。

这样即使 Q1--Q8 中若干问题老师后续回答不同，
最小闭环的结果\textbf{都能作为讨论基础}。
\end{keybox}

% ============ Section 6: Summary ============
\section{总结}

\begin{keybox}[title=\textbf{本文核心}]
\begin{enumerate}[leftmargin=1.5em, itemsep=2pt]
\item Lagendijk 2026 论文的核心是 BCH $t \le 4$ 的 \textbf{Direct root finding} 硬件译码器，
      对 $t=2$ 相对 Conventional 有 62\% 时延压缩（3 cyc vs 8 cyc）。
\item 老师本次任务是把 Direct BCH 硬判 + RS-BM 硬判\textbf{级联}，在 PAM4-AWGN 信道下
      比较与纯 RS-BM 的 BER 和时延，KPI 是级联时延 $\le$ RS-BM + 10\%。
\item 相对第一阶段（软判决），本次\textbf{只换内码 BCH}（LLOSD $\to$ Direct t=2），
      $\sim$70\% 基础设施可复用。
\item 有 8 个关键技术问题需要对齐（\S4），其中 Q1--Q3 会\textbf{根本决定路线}。
\item 建议先做 \textbf{$n=255$ + Direct + Conv + RS-BM} 最小闭环 (预计 3--4 天)，
      拿到中间结果后再横向扩展。
\end{enumerate}
\end{keybox}

\vspace{6pt}
\hrule
\vspace{4pt}
\noindent\footnotesize{
\textbf{参考文献}\\
{[1]} J.\ Lagendijk et al., ``High-throughput Low-latency Hardware Implementation of BCH Decoders,''
arxiv 2606.17837v1, EPFL + Eindhoven Univ.\ of Technology, Jun 2026.\\
{[2]} L.\ Yang et al., ``Efficient OSD of BCH Codes Without GE,''
{\it IEEE Trans.\ Inf.\ Theory}, vol.\ 71, no.\ 11, pp.\ 8294--8311, Nov 2025.\\
\textbf{第一阶段代码}: \url{https://github.com/a-yeyang/efficient-osd-bch-no-ge/tree/main/rs_bch_cascade}\\
\textbf{文档}: 2026-07-23 · 陈诗阳
}

\end{document}
